\clearpage
\fancyhead[L]{\small TOBIAS MAYER. ASTRONOMICAL OBSERVATIONS.}
\hypertarget{native-649-lower}{}
\begin{center}\small Göttingische gelehrte Anzeigen. Stück 160. Seite 1585..1588. 1827 October 6.\\[3pt]\rule{.72\linewidth}{.3pt}\end{center}

\emph{Astronomical observations made at Göttingen from \textup{1756} to \textup{1761} by \textup{\textsc{Tobias Mayer}}. In two parts. Published by order of the commissioners of latitude.} London 1826. Die erste Abtheilung 92 Seiten, die zweite 62 Seiten in Folio.

Nach siebenzig Jahren ist den Beobachtungen des grossen Göttinger Astronomen \textsc{Tobias Mayer} noch die Auszeichnung zu Theil geworden, auf Kosten einer ausländischen Behörde, welche die grossen ihr zu Gebote stehenden Mittel gern zum Vortheil der Astronomie anwendet, gedruckt zu werden. Die Handschrift der Beobachtungen von den Jahren 1757..1761 war schon vor langer Zeit in den Besitz des Herrn von \textsc{Zach} gekommen, und von diesem der englischen Commission für die Meereslänge mitgetheilt: die Handschrift der Beobachtungen des ersten Jahres befand sich noch in dem Archiv der hiesigen königl. Societät und wurde zu dem beabsichtigten Abdrucke hergeliehen.

Der bei weitem grösste Theil der Beobachtungen ist an dem \textsc{Bird}'schen Mauerquadranten gemacht. \textsc{Mayer} hatte sein Hauptaugenmerk auf die Verfertigung des Verzeichnisses der Zodiakalsterne gerichtet, und die Beobachtungen dieser Sterne sind daher auch der Hauptinhalt des Tagebuchs. Sie fangen mit dem 8. Februar 1756 an; zuerst bis zum 23. Mai finden sich blos die Durchgangszeiten, und zwar schon auf den mittlern Faden reducirt; vom 23. Mai an hingegen erscheinen die Beobachtungen vollständig im Originale. Die Beobachtungen des Jahrs 1756 füllen die ganze erste Abtheilung aus; der Jahrgang 1757 nimmt in der zweiten Abtheilung 35 Seiten ein; dann 6 Seiten bis in die Mitte des Jahrs 1758. Hiermit hören eigentlich \textsc{Mayers} Beobachtungen am Mauerquadranten fast ganz auf; am 5. April 1759 übergab er sein Verzeichniss der Zodiakalsterne der Königlichen Societät. Aus den ersten Tagen des Mais 1759 Beobachtungen des \textsc{Halley}'schen Cometen an einem parallactisch aufgestellten Fernrohr. Aus dem Juni und Juli 1760 Beobachtungen am Mauerquadranten von \textsc{Mayers} würdigem Schüler \textsc{Niebuhr}; endlich vom 6. Juni 1761 die Beobachtung des Durchganges der Venus durch die Sonne. Neben den Beobachtungen der Zodiakalsterne finden wir, in den ersten Jahren, noch sehr viele Beobachtungen der Sonne, mehrere von der Venus und dem Monde, und einige wenige vom Mars, Jupiter und Saturn, auch verschiedene Sternbedeckungen.

Man sieht aus diesem Inhalt, dass eine neue von Grund aus gemachte Bearbeitung dieser Beobachtungen keine so umfassende Ausbeute liefern könnte, wie die gleichzeitigen Beobachtungen von \textsc{Bradley} durch \textsc{Bessels} vortreffliche Bearbeitung gegeben haben. \textsc{Mayer} hatte sich ein beschränktes \SourceLeaf{650}Hauptziel vorgesetzt, die Anfertigung eines Zodiakalsternenverzeichnisses, und dieses hat er selbst auf seine Beobachtungen gegründet. Auch darf nicht übersehen werden, dass die Dürftigkeit von \textsc{Mayer}'s Hülfsmitteln, für die Reduction der Beobachtungen als absoluter, grosse Hindernisse in den Weg legt, und dass uns jetzt die Kenntniss von einigen Nebenhülfsmitteln, die er noch benutzte, abgeht. \textsc{Mayer} hatte kein Mittagsfernrohr; die Reduction der am Mauerquadranten beobachteten geraden Aufsteigungen setzt die Kenntniss der unregelmässigen Abweichungen des Instruments vom Meridian voraus, wozu \textsc{Mayer} auch correspondirende Sonnenhöhen an einem kleinen beweglichen Quadranten benutzte; allein nur von zwei Tagen finden wir diese aufbewahrt. Aehnliche Schwierigkeiten finden sich bei den Zenithdistanzen. \textsc{Mayer} hatte keinen Zenithsector. Eine zuverlässige Kenntniss des Collimationsfehlers konnte daher nur durch Umhängen erlangt werden, und dies Umhängen ist nur Einmal geschehen. Alle vor dem Umhängen beobachteten Zenith'distanzen verlieren noch durch die von \textsc{Mayer} selbst gemachte Bemerkung an Zuverlässigkeit, dass der Centralzapfen lose sitzend befunden wurde. Es liessen sich noch mehrere andere Bemerkungen beifügen, wozu aber hier nicht der Ort ist.

Dagegen werden noch oft einzelne Fälle eintreten können, wo es den Astronomen sehr willkommen sein wird, die Originalbeobachtungen benutzen zu können, wenn auch nur als Differentialbeobachtungen. Bekanntlich hat schon einmal die Handschrift den grossen Nutzen geleistet, eine frühe Beobachtung des Uranus zu liefern, den \textsc{Mayer} am 25. Sept. 1756 unbewussterweise als Fixstern beobachtete.

In der Einleitung sind noch zwei Aufsätze der Opera inedita wieder abgedruckt, nemlich \textsc{Mayer}'s Vorlesung Observationes astronomicae quadrante murali habitae, und \textsc{Lichtenberg}'s Bemerkungen über das von \textsc{Mayer} gebrauchte Thermometer; und am Schluss findet sich ein vom Hrn. Prof. \textsc{Bessel} eingeliefertes Druckfehlerverzeichniss zu \textsc{Bradley}'s Beobachtungen.
